\section{Manual (Aplicación socios) - Apartado de reservas}

Para utilizar los distintos espacios del club, el socio debe ingresar a la sección de reservar instalación. Allí podrá visualizar todas las opciones disponibles, notando que algunas tienen un límite de personas y un precio asignado, mientras que otras son de uso gratuito (como, por ejemplo, la parrilla).

\subsection{Reserva de instalaciones sin costo}
El proceso para reservar un espacio gratuito es directo:
\begin{itemize}
    \item Se debe seleccionar la instalación y elegir uno de los turnos disponibles.
    \item Al confirmar la operación, el sistema redirigirá a la sección de mis reservas.
    \item Como el espacio no tiene un precio asociado, la reserva quedará aprobada y su estado será directamente confirmada.
\end{itemize}

\subsection{Reserva de instalaciones aranceladas y pago}
Cuando se trata de un espacio con costo (por ejemplo, una cancha de tenis), el flujo requiere un paso adicional para efectivizar el turno:
\begin{itemize}
    \item Se selecciona la instalación, el día y el horario deseado.
    \item Al confirmar, la solicitud aparecerá en la sección de mis reservas con el estado pendiente, ya que es necesario realizar el abono.
    \item Para confirmarla, se debe seleccionar la opción de pago, la cual derivará al usuario a la pasarela de Mercado Pago (un proceso idéntico al pago de la cuota social).
    \item Una vez que se elige el medio de pago y la transacción es aceptada, el estado del turno se actualizará automáticamente a confirmada.
\end{itemize}
